\subsection{VALUED DIVISION RINGS}

DEFINITION 2. An absolute value on a division ring K is a mapping \(x \to |x|\) of K into \(\mathbf{R}_{+}\) which satisfies the following conditions:

\((\mathbf{VM}_{\mathbf{I}})\) \(|x| = 0\) if and only if \(x = 0\) .

\((\mathbf{VM}_{\mathbf{II}})\quad |xy| = |x|.|y|\) for all \(x,y\) in \(\mathbf{K}\) .

\((\mathbf{VM}_{\mathbf{III}})\quad |x + y|\leqslant |x| + |y|\) for all \(x,y\) in \(\mathbf{K}\) .

By \((\mathbf{VM}_{\mathbf{II}})\) we have \(|x| = |\mathrm{I}|.|x|\) , and since by \((\mathbf{VM}_{\mathbf{I}})\) there is at least one \(x\) such that \(|x| \neq 0\) , we have \(|\mathrm{I}| = \mathrm{I}\) ; it follows that \(\mathrm{I} = |- \mathrm{I}|^2\) , hence \(|- \mathrm{I}| = \mathrm{I}\) and consequently

\[
| - x | = | - \mathrm{I} | | x | = | x |;
\]

therefore \(|x - y| \leqslant |x| + |y|\) for all \(x, y\) in \(\mathbf{K}\) . We can therefore say that \(d(x, y) = |x - y|\) is an invariant metric on the additive group \(\mathbf{K}\) , and that the mapping \(x \to |x|\) is a homomorphism of the multiplicative group \(K^{*}\) of non-zero elements of K into the multiplicative group \(R_{+}^{*}\) of real numbers \textgreater0.

The invariant metric \(|x - y|\) defines a metric space topology on K, compatible with its additive group structure (no. 1, Proposition 3); but, moreover, this topology is compatible with the division ring structure of K. For the continuity of xy on \(K \times K\) follows from the relation

\[
x y - x _ {0} y _ {0} = (x - x _ {0}) (y - y _ {0}) + (x - x _ {0}) y _ {0} + x _ {0} (y - y _ {0}),
\]

which gives

\[
\left| x y - x _ {0} y _ {0} \right| \leqslant \left| x - x _ {0} \right|. \left| y - y _ {0} \right| + \left| x _ {0} \right|. \left| y - y _ {0} \right| + \left| y _ {0} \right|. \left| x - x _ {0} \right|.
\]

Likewise, the continuity of \(x^{-1}\) at every point \(x_{0} \neq 0\) follows from the identity \(x^{-1} - x_{0}^{-1} = x^{-1} (x_{0} - x)x_{0}^{-1}\) , which gives, by \((\mathrm{VM}_{\mathrm{II}})\) ,

\[
\left| x ^ {- 1} - x _ {0} ^ {- 1} \right| = \frac {\left| x - x _ {0} \right|}{\left| x _ {0} \right| . \left| x \right|};
\]

now if \(\varepsilon > 0\) is such that \(\varepsilon < |x_0|\) , the relation \(|x - x_0| \leqslant \varepsilon\) implies \(|x| \geqslant |x_0| - \varepsilon\) , whence \(|x^{-1} - x_{0}^{-1}| \leqslant \varepsilon / |x_0| (|x_0| - \varepsilon)\) ; and this establishes the continuity of \(x^{-1}\) at the point \(x_0\) .

DEFINITION 3. A valued division ring is a division ring K endowed with the structure defined by a given absolute value on K.

A valued division ring will always be considered as endowed with the topology defined by its absolute value, which makes it a topological division ring. If \(K_{0}\) is a division subring of a valued division ring K, the restriction to \(K_{0}\) of the absolute value on K is an absolute value on \(K_{0}\) , which defines on \(K_{0}\) the topology induced by the topology of K.

Examples. 1) Let K be an arbitrary division ring. For each \(x \in \mathbf{K}\) , put \(|x| = \mathbf{i}\) if \(x \neq 0\) , and \(|0| = 0\) . The mapping \(x \to |x|\) so defined is an absolute value on K, called the improper absolute value. The topology defined by an absolute value \(|x|\) on a division ring K is discrete if and only if \(|x|\) is the improper absolute value. This condition is clearly sufficient; conversely, if the topology of K is discrete, \(|x|\) can take no value \(\alpha > 0\) other than \(\mathbf{i}\) ; for if we had \(|x_0| = \alpha < \mathbf{i}\) , the sequence \((x_0^n)\) would consist of non-zero terms and would converge to \(\mathbf{o}\) ; to deal with the case \(\alpha > \mathbf{i}\) , consider \(x_0^{-1}\) in place of \(x_0\) .

\begin{enumerate}
\def\labelenumi{\arabic{enumi})}
\setcounter{enumi}{1}
\item
  The absolute value of a real number (Chapter IV, § 1, no. 6) satisfies axioms \((\mathrm{VM}_{\mathrm{I}})\) , \((\mathrm{VM}_{\mathrm{II}})\) and \((\mathrm{VM}_{\mathrm{III}})\) , and the topology it defines on the field \(\mathbf{R}\) is the topology of the real line. On the field \(\mathbf{C}\) of complex numbers (identified with \(\mathbf{R}^2\) ) {[}resp. the division ring \(\mathbf{H}\) of quaternions (identified with \(\mathbf{R}^4\) ){]} the Euclidean norm is again an absolute value and defines the topology of the field C (resp. the division ring H) (Chapter VIII, § 1, nos. 2 and 4).
\item
  On a division ring K, a real valuation is a function v defined on \(K^{*}\) with values in R which satisfies the following conditions: a) if \(x \in K^{*}\) and \(y \in K^{*}\) , then \(v(xy) = v(x) + v(y)\) ; b) if in addition \(x + y \neq 0\) , then \(v(x + y) \geqslant \inf(v(x), v(y))\) . If a is any real number \textgreater1, we can then define an absolute value on K by putting \(|x| = a^{-v(x)}\) for \(x \neq 0\) , and \(|0| = 0\) . For the relation \(v(xy) = v(x) + v(y)\) for \(x \neq 0\) and \(y \neq 0\) implies the relation \(|xy| = |x|.|y|\) for these values of x and y, and this relation is trivially true if one of x, y is zero; likewise, from the relation \(v(x + y) \geqslant \inf(v(x), v(y))\) for \(x \neq 0\) , \(y \neq 0\) and \(x + y \neq 0\) we deduce \(|x + y| \leqslant \sup(|x|, |y|) \leqslant |x| + |y|\) , and these inequalities are still satisfied if one of x, y, \(x + y\) is zero. In particular, if \(v_{p}(x)\) is the p-adic valuation on the field Q of rational numbers (the exponent of p in the decomposition of x into a product of prime factors), then the corresponding absolute value \(|x|_{p} = p^{-v_{p}(x)}\) is called the p-adic absolute value on the field Q (cf.~Chapter III, §6, Exercise 23).
\end{enumerate}

Remark. If \(x\) is a root of unity in a valued division ring, then \(|x| = 1\) , for \(x^n = 1\) implies \(|x|^n = 1\) and hence \(|x| = 1\) . In particular, the only absolute value on a finite field is the improper absolute value, since every element \(\neq 0\) of such a field is a root of unity.

DEFINITION 4. Two absolute values on a division ring K are said to be equivalent if they define the same topology on K.

PROPOSITION 5. Two absolute values \(|x|_1, |x|_2\) on a division ring \(K\) , neither of which is the improper absolute value, are equivalent if and only if the relation \(|x|_1 < i\) implies \(|x|_2 < i\) . There exists then a real number \(\rho > 0\) such that \(|x|_2 = |x|_1^p\) for all \(x \in K\) .

The condition is necessary, for the set of all \(x \in \mathbf{K}\) such that \(|x|_1 < 1\) is the same as the set of all \(x\) such that, with respect to the topology defined by the absolute value \(|x|_1\) , \(\lim x^n = 0\) .

Suppose conversely that \(|x|_1 < \mathrm{i} \Longrightarrow |x|_2 < \mathrm{i}\) . Then \(|x|_1 > \mathrm{i} \Longrightarrow |x|_2 > \mathrm{i}\) , because \(|x^{-1}|_1 < \mathrm{i}\) and therefore \(|x^{-1}|_2 < \mathrm{i}\) . Since by hypothesis the absolute value \(|x|_1\) is not improper, there exists \(x_0 \in \mathbf{K}\) such that \(|x_0|_1 > \mathrm{i}\) . Let \(a = |x_0|_1\) , \(b = |x_0|_2\) and let \(\rho = \log b / \log a > 0\) . Let \(x \in \mathbf{K}^*\) and put \(|x|_1 = |x_0|^{\gamma}\) . If \(m\) and \(n\) are integers such that \(n > 0\) and \(m/n > \gamma\) , then \(|x|_1 < |x_0|_1^{m/n}\) , and therefore \(|x^n x_0^{-m}|_1 < \mathrm{i}\) ; hence \(|x^n x_0^{-m}|_2 < \mathrm{i}\) , \(|x|_2 < |x_0|_2^{m/n}\) . Similarly, if \(m/n < \gamma\) , we see that \(|x|_2 > |x_0|_2^{m/n}\) ; it follows therefore that \(|x|_2 = |x_0|_2^\gamma\) ; in other words

\[
\log | x | _ {2} = \gamma \log b = \gamma \rho \log a = \rho \log | x | _ {1},
\]

i.e., \(|x|_2 = |x|_1^p\) . It is now clear that the neighbourhoods of zero for the topologies defined on \(\mathbf{K}\) by \(|x|_1\) and \(|x|_2\) are identical.

Conversely, if \(|x|\) is any absolute value on \(K\) , the function \(|x|^{\circ}\) is an absolute value on \(K\) (equivalent to \(|x|\) ) for all \(\rho\) such that \(0 < \rho \leqslant 1\) . We have only to verify the inequality \(|x + y|^{\circ} \leqslant |x|^{\circ} + |y|^{\circ}\) ; and since \(|x + y|^{\circ} \leqslant (|x| + |y|)^{\circ}\) it is enough to show that, if \(a > 0\) and \(b > 0\) , we have \((a + b)^{\circ} \leqslant a^{\circ} + b^{\circ}\) for any \(\rho\) such that \(0 < \rho \leqslant 1\) . If we put \(c = a / (a + b)\) and \(d = b / (a + b)\) , we have \(c + d = 1\) , and the inequality to be proved is \(c^{\circ} + d^{\circ} \geqslant 1\) ; but this follows immediately from the relations \(c^{\circ} \geqslant c\) and \(d^{\circ} \geqslant d\) , which are valid since \(0 < c \leqslant 1\) , \(0 < d \leqslant 1\) , \(0 < \rho \leqslant 1\) .

Hence the set of values of r \textgreater{} 0 such that \(|x|^{r}\) is an absolute value is a finite or infinite interval of R with left-hand end-point o; if it is finite, it is evidently closed; for if we have \(|x + y|^{r} \leqslant |x|^{r} + |y|^{r}\) for any x, y in K and all r such that \(0 < r < r_{0}\) , then by continuity the inequality is still valid for \(r = r_{0}\) . If \(|x|^{r}\) is an absolute value for all r \textgreater{} 0, then we have

\[
| x + y | \leqslant (| x | ^ {r} + | y | ^ {r}) ^ {1 / r}
\]

for all \(x\) and \(y\) in \(\mathbf{K}\) and all \(r > 0\) . Now, if \(a, b\) are two real numbers \(\geqslant 0\) , we have \(\lim_{r \to \infty} (a^r + b^r)^{1/c} = \sup (a, b)\) ; for, supposing for example that \(a \geqslant b\) , we have \(a \leqslant (a^r + b^r)^{1/r} \leqslant 2^{1/r} a\) , and the result follows by letting \(r \to +\infty\) .

Thus, if \(|x|^r\) is an absolute value for all \(r > 0\) , we have

\[
| x + y | \leqslant \sup (| x |, | y |)
\]

which can be expressed by saying that \(v(x) = -\log |x| (x \neq 0)\) is a valuation on K.

Remark. The proof of Proposition 5 shows that, if the topology defined by \(|x|_2\) is coarser than that defined by \(|x|_1\) , and if \(|x|_1\) is not improper, then \(|x|_1\) and \(|x|_2\) are equivalent, for the relation \(|x|_2 < 1\) then implies \(|x|_2 < 1\) . Thus the topologies defined by two absolute values on \(K\) , neither of which is improper, cannot be comparable without being identical.

PROPOSITION 6. The completion \(\hat{\mathbf{K}}\) of a division ring \(\mathbf{K}\) endowed with an absolute value \(|x|\) is a division ring, and the function \(|x|\) can be extended by continuity to an absolute value on \(\hat{\mathbf{K}}\) , which defines the topology of \(\hat{\mathbf{K}}\) .

Let \(\mathfrak{F}\) be a Cauchy filter on K (with respect to the additive uniformity) which does not have o as a cluster point; to show that \(\hat{\mathbf{K}}\) is a division ring it is enough to establish that the image of \(\mathfrak{F}\) under the mapping \(x \to x^{-1}\) is a Cauchy filter base (Chapter III, § 6, no. 8, Proposition 7). Now, by hypothesis there exists a real number \(\alpha > 0\) and a set \(\mathrm{A} \in \mathfrak{F}\) such that \(|x| \geqslant \alpha\) for all \(x \in \mathrm{A}\) ; on the other hand, for each \(\varepsilon > 0\) there exists a set \(\mathrm{B} \in \mathfrak{F}\) such that \(\mathrm{B} \subset \mathrm{A}\) and \(|x - y| \leqslant \varepsilon\) for all \(x \in \mathrm{B}\)

and \(y\in \mathbf{B}\) ; hence

\[
\left| x ^ {- 1} - y ^ {- 1} \right| = \frac {\left| x - y \right|}{\left| x \right| . \left| y \right|} \leqslant \frac {\varepsilon}{\alpha^ {2}},
\]

and the first part of the proposition follows. The invariant metric \(|x - y| = d(x, y)\) extends by continuity to a metric on \(\hat{\mathbf{K}}\) (§ 2, no. 1, Proposition 1) which defines the topology of \(\hat{\mathbf{K}}\) and is invariant by the principle of extension of identities; we continue to denote this invariant metric by \(d(x, y)\) . If we put \(|x| = d(0, x)\) for \(x \in \hat{\mathbf{K}}\) , it is clear that \(|x|\) is the extension by continuity of the function \(|x|\) on \(\mathbf{K}\) and is therefore an absolute value on \(\hat{\mathbf{K}}\) by the principle of extension of identities.

